\scanpage{9}
\mspage{8}

\uline{Lieux disjoints} \inserted{$C,D$} sur une ficelle \inserted{(relation symétrique)}:

\deleted{ou bien $C$}
Si $C$ ou $D$ est une extrémité, \uline{disjoint} signifie que $C\ne D$. Si ni $C$ ni $D$ n'est une extrémité, alors disjoint signifie que
\deleted{\unsure{Si}} $C\ne D$, \uline{et} que $D$ est un lieu de $F_{A,C}$ ou de $F_{B,C}$ ($\Rightarrow C$ est un lieu de $F_{B,D}$ ou de $F_{A,D}$).

\uline{Ficelle orientée} Ficelle avec extrémité fixée (``origine).

\deleted{\unsure{A}} La catégorie des ficelles \inserted{orientées} munies de deux lieux internes \uline{disjoints}, est équivalente à celle des \deleted{systèmes} \deleted{triplets}\inserted{quintuplets} $(F_1,\unsure{F_1|F_2},F_3,h_{21},h_{32})$, où $F_1,F_2,F_3$ sont des ficelles orientées,
\[
h_{21}:\operatorname{ex}(F_1)\simeq\unsure{\operatorname{or}}(F_2),\quad
h_{32}:\operatorname{ex}(F_2)\simeq\operatorname{or}(F_3).
\]

$F=(X,(A,B))$ ficelle ordonnée. Soit $C$ un lieu, on pose $\unsure{\mathcal F_C|F_C}={}\text{ens des lieux de }F_{A,C}$.
\margin{si $C\ne A,B$; $=\{A\}$ si $C=A$, $=\unsure{F|\mathcal L}$ si $C=B$.}
On écrit $C\le D$ ssi $F_C\subset F_D$. Cela signifie aussi que $C\in F_D$.

\rotatebox{90}{\reflectbox{\unsure{EXTRA STRONG}}}
\rotatebox{90}{\reflectbox{\unsure{EXTRA STRONG}}}
